\section{Growing Windows and Reduction to a Base Length}
\label{sec:growing-windows}

The statements below distinguish two answers to the growing-window question.
Small growing updates still encounter the balanced bottleneck, whereas a
specified ideal heat bath with windows of order $N\log N$ mixes rapidly at
every feasible length.  The latter counts completed exact conditional draws;
it does not supply a polynomial-time implementation of those draws over all
transpositions.  All products act rightmost first and windows do not wrap.

\begin{corollary}[Obstruction for growing update budgets]
\label{cor:growing-obstruction}
Use the parameters $N_m,K$ and balanced event $A_K$ of
Theorem~\ref{thm:balanced-bottleneck}.  Let $b=b_m\ge1$ be integer and
$b=o(m/(\log m)^3)$.  Every kernel $L_m$ with uniform stationary law on
$\Omega(N_m,m)$ that changes at most $b$ gate positions per step satisfies
\[
 \frac{Q_{L_m}(A_K,A_K^c)}{\min\{\pi(A_K),\pi(A_K^c)\}}
 \le C N_m4^{-s_b},\qquad
 s_b=\left\lfloor\frac{K-b/2}{b+1}\right\rfloor+1.
\]
Here $s_b/\log m\to\infty$.  At time $t_m=\lfloor2^{s_b}\rfloor$, a
$1-o(1)$ fraction of starting words have TV distance at least $1/2-o(1)$.
In particular, mixing is superpolynomial in $N_m,m$.  Neither consecutive
positions nor reversibility is required.
\end{corollary}
\begin{proof}
Equation (CB.6) gives the cut bound, with an absolute constant because its
counting estimate is uniform in $b$.  Since
$K\gtrsim m/(\log m)^2$ and $b=o(m/(\log m)^3)$,
$s_b\asymp K/(b+1)=\omega(\log m)$.  Thus
$t_m Q_{L_m}(A_K,A_K^c)\le C N_m2^{-s_b}=o(1)$.
The stationary crossing and Markov-inequality argument in
Theorem~\ref{thm:balanced-bottleneck} proves the individual-start claim.
The selected $N_m$ is polynomially bounded in $m$.
\end{proof}

For the positive result define the endpoint fiber
$\Omega_g(N,m)=\{w\in T_N^m:\prod_iw_i=g\}$, with its uniform law.
Let $\mu_{N,r}$ be the product distribution of $r$ independent uniform
transpositions.  By (CB.7), the choice
\[
 B=\lceil8N\log(N+1)\rceil
 \quad\hbox{gives}\quad
 \left|\frac{N!}{2}\mu_{N,r}(x)-1\right|\le\eta:=\frac1{16}
 \quad(r\ge B,\ \operatorname{sgn}(x)=(-1)^r).       \tag{GW.1}
\]
We take $N\ge3$; the $N=2$ fibers are singletons.

\begin{theorem}[Rapid mixing of an ideal growing-window heat bath]
\label{thm:growing-window-mixing}
For every nonempty $\Omega_g(N,m)$, set $q=\max\{1,\lfloor m/B\rfloor\}$.
For $q\ge2$, partition the word into $q$ blocks: the first $q-1$ have
length $B$ and the last has length in $[B,2B-1]$.  With probability $1/2$
hold; otherwise choose one of the $d=q-1$ adjacent block pairs uniformly
and refresh its union by an exact uniform same-product draw.  Call this
kernel $\mathcal B$.  For $q=1$, refresh the whole word with probability
$1/2$.  Then every move changes at most $\min\{m,3B\}$ positions and
\[
 \operatorname{gap}(\mathcal B)\ge\frac1{16(q-1)},\qquad
 \Phi(\mathcal B)\ge\frac1{32(q-1)}\quad(q\ge2),       \tag{GW.2}
\]
\[
 t_{\rm mix}(\mathcal B,\epsilon)
 \le \lceil4d\log(2d/\epsilon)\rceil
      +\lceil2d\log(2q/\epsilon)\rceil.              \tag{GW.3}
\]
For $q=1$ the gap is $1/2$ and $t_{\rm mix}\le\lceil\log_2(1/\epsilon)\rceil$.
Singleton gaps and conductances may be set to one.  These bounds are
uniform in the endpoint $g$ and count ideal block updates.
\end{theorem}
\begin{proof}
Let $x_i$ be the prefix product at block boundary $i$, with $x_0=e,x_q=g$.
Their joint density is proportional to
$\prod_i\mu_{N,L_i}(x_{i-1}^{-1}x_i)$.
Given these boundaries, block interiors are independent uniform bridges.
Refreshing blocks $i,i+1$ first heat-bath updates $x_i$ and then draws
both interiors conditional on the new boundaries.

By (GW.1), every conditional law of $x_i$ dominates
$a=((1-\eta)/(1+\eta))^2$ times the uniform law on its prescribed parity
class.  Any two such conditionals admit a coupling with disagreement
probability at most $\beta=1-a\le64/289$.  For boundary vectors differing
only at $j$, updating $j$ erases the disagreement; only $j-1,j+1$ can
create another.  Half-laziness therefore gives expected Hamming distance
at most $1-(1-2\beta)/(2d)\le1-1/(4d)$.
Every parity-compatible boundary vector is feasible by (GW.1), so
telescoping along Hamming paths extends the contraction to any pair of
vectors.  After $t$ steps their mismatch probability is at most
$d e^{-t/(4d)}$.

Once the boundaries agree, use identical subsequent draws.  Each block
is refreshed with probability at least $1/(2d)$ per step, so a further
$t$ steps leave any interior uncoupled with probability at most
$q e^{-t/(2d)}$.  This proves (GW.3).  Splitting a long time into two
equal phases gives exponential decay at rate $1/(8d)$.  Applying that
bound to an eigenfunction of this finite reversible half-lazy kernel
gives gap at least $1-e^{-1/(8d)}\ge1/(16d)$.
The indicator variational inequality gives $\Phi\ge\operatorname{gap}/2$.
\end{proof}

\begin{theorem}[All-endpoint base-length reduction]
\label{thm:base-length-reduction}
Let $P_{N,\ell,g}$ be the half-lazy adjacent pair heat bath on
$\Omega_g(N,\ell)$, with the same normalization as $P_{N,m}$.
Define
\[
 \phi_* =\min_{\substack{2B\le\ell\le4B-2\\
                   \operatorname{sgn}(h)=(-1)^\ell}}
                    \Phi(P_{N,\ell,h}),\qquad
 \gamma_* =\min_{\substack{2B\le\ell\le4B-2\\
                   \operatorname{sgn}(h)=(-1)^\ell}}
                    \operatorname{gap}(P_{N,\ell,h}).
\]
For every $m\ge2B$ and feasible $g$,
\[
 \Phi(P_{N,m,g})\ge\frac{2B-1}{32(m-1)}\phi_*,\qquad
 \operatorname{gap}(P_{N,m,g})
       \ge\frac{2B-1}{16(m-1)}\gamma_*.              \tag{GW.4}
\]
The implication is conditional as a rapid-mixing statement: no
inverse-polynomial lower bound for $\phi_*$ or $\gamma_*$ is supplied.
\end{theorem}
\begin{proof}
Use the block partition from Theorem~\ref{thm:growing-window-mixing},
writing $I_i$ for a neighboring union.  Its length $\ell_i$ lies in
$[2B,4B-2]$.  For $f=1_A$, the local conditional cut with mass $p$
satisfies $Q_i\ge\phi_*\min(p,1-p)\ge\phi_*p(1-p)$.
The ideal-chain Poincare bound gives
\[
 \operatorname{Var}_\pi f
 \le8\sum_i\mathbb E\operatorname{Var}(f\mid I_i^c)
 \le\frac8{\phi_*}\sum_i\mathbb E Q_i
 \le\frac{16(m-1)}{(2B-1)\phi_*}Q_P(A,A^c).
\]
For the last inequality, each adjacent bond is in at most two unions,
and its local coefficient $1/[2(\ell_i-1)]$ is at most
$(m-1)/(2B-1)$ times its global coefficient.  If $\pi(A)\le1/2$,
$\operatorname{Var}1_A\ge\pi(A)/2$, proving the conductance inequality.
For general $f$, use the local Poincare inequality with $\gamma_*$ in
place of the conditional cut bound; this gives the gap inequality.
\end{proof}

The two useful directions are now explicit.  A uniform inverse-polynomial
bound over \emph{all endpoints} at lengths in $[2B,4B-2]$ implies polynomial
mixing at every longer polynomial length, since
$\log|\Omega_g(N,m)|\le m\log\binom N2$ and Cheeger's inequality converts
conductance to gap.  Conversely, a superpolynomially small conductance at
some polynomial length $m\ge2B$, including $m\asymp N^2$, forces
superpolynomially small conductance in at least one of these base fibers.
This is a contrapositive, not a claim that every short bottleneck persists
at every longer length.  It neither controls $m<2B$ nor replaces all
endpoints by $h=e$.  The intermediate bottleneck family has $m\ll N\log N$
and therefore does not resolve the base-length question.

The ideal-chain theorem settles a mixing question for a specified growing
window rule, not for every growing window size or scan rule.  In particular,
when $m<2B$ it uses a whole-word refresh.  Exact large-block sampling cost
remains a separate obligation.  Odd identity lengths are empty; $m=0$ is
the empty singleton and $m=2$ uses the whole-word rule.  None of these
statements turns a bound polynomial in $N,m$ into a bound polynomial in
the wire count $\log_2N$.
